% AAAI 2027 VA-ATACOM Supplementary Material (standalone, separate PDF)
\documentclass[letterpaper]{article}
\usepackage[submission]{aaai2027}
% NOTE: AAAI sections are unnumbered (secnumdepth=0), so \ref to a section
% resolves to empty. Cross-references to the main paper are written as plain
% text ('Table 1 of the main paper' etc.) for reviewer clarity.
\usepackage[hyphens]{url}
\usepackage{graphicx}
\urlstyle{rm}
\def\UrlFont{\rm}
\usepackage{natbib}
\usepackage{caption}
\usepackage{amsmath,amssymb,amsfonts}
\usepackage{booktabs}
\usepackage{multirow}
\usepackage{xcolor}
\frenchspacing

\pdfinfo{
/TemplateVersion (2027.1)
}

\setcounter{secnumdepth}{0}

\title{VA-ATACOM Supplementary Material:\\The DT-margin Heuristic Baseline}

\author{
    % Authors hidden by [submission] option
    Anonymous Submission
}
\affiliations{
    Anonymous Institution
}

\begin{document}
\maketitle

\section{The DT-margin Heuristic Baseline}
\label{app:dtmargin}

VA-ATACOM's velocity augmentation has a heuristic predecessor we call the \textbf{DT-margin} filter --- the ``DT-margin (heuristic)'' row of Table~1 in the main paper. It inflates the keepout radius linearly with speed, $r_\text{eff}(\boldsymbol{v}) = r_\text{base}(1+\alpha\|\boldsymbol{v}\|)$, projects out the inward radial component, and adds the same sim-BRT reward shaping. As shown in the main paper's design study, this linear margin is the \emph{first-order} approximation of VA-ATACOM's quadratic braking constraint; under fair hazard-radius matching it reaches 13/15 GO --- two Push seeds short of VA-ATACOM's 15/15, illustrating that the linearisation almost captures the right behaviour but leaves a residual gap that only the exact quadratic constraint closes. We characterise it here for completeness --- its two tunable knobs ($\alpha$ and the BRT horizon $h$) and their interaction with the control rate. These findings motivated, and are superseded by, the parameter-free velocity-augmented manifold of the main paper.

\subsection{Ablation of the two components}

\begin{table}[t]
\centering
\caption{DT-margin ablation on Goal (5 seeds).}
\label{tab:ablation}
\setlength{\tabcolsep}{4pt}
\begin{tabular}{lcccc}
\toprule
Config & $\alpha$ & $h$ & Mean $\bar{C}$ & GO/5 \\
\midrule
Velocity-adaptive only & 0.3 & 0 & 2.90 & 5/5 \\
BRT-only & 0.0 & 3 & 2.35 & 4/5 \\
\textbf{DT-margin (full)} & 0.3 & 3 & \textbf{0.96} & \textbf{5/5} \\
\bottomrule
\end{tabular}
\end{table}

% Inline figure with reduced width to fit 2-page layout.
\begin{figure}[t]
  \centering
  \includegraphics[width=0.85\columnwidth]{figures/fig_v_c_ablation_bar/fig_v_c_ablation_bar.pdf}
  \caption{DT-margin ablation: both components are necessary. Mean episode cost $\bar C$ on SafetyPointGoal1-v0 (5 seeds $\times$ 200K env steps), bars are $\pm 1$ std. VAM-only ($\alpha{=}0.3, h{=}0$) addresses M1 chord excursion only; BRT-only ($\alpha{=}0$, $h{=}3$) anticipates M2 drift but one seed collapses via M1; DT-margin ($\alpha{=}0.3$, $h{=}3$) combines both for lowest mean (0.96) and 5/5 GO.}
  \label{fig:v_c_ablation_bar}
\end{figure}

Table~\ref{tab:ablation} (Fig.~\ref{fig:v_c_ablation_bar}) shows both components contribute: velocity-adaptive margin alone reaches 5/5 GO at higher variance (2.90); BRT alone 4/5 GO (one M1 collapse at $C{=}7.98$); the full heuristic 0.96 cost / 100\% GO.

\subsection{$\Delta t$ sensitivity}
\label{sec:dt-sweep}

\begin{table}[t]
\centering
\caption{DT-margin cost vs.\ time step ($\bar C \pm \sigma$, 5 seeds, 50K steps).}
\label{tab:dt}
\setlength{\tabcolsep}{4pt}
\begin{tabular}{lccc}
\toprule
$\Delta t$ & ATACOM & DT-margin ($h{=}3$) & Note \\
\midrule
0.01s  & 0.00  & 3.02 & ATACOM OK \\
0.02s  & 0.00  & 0.00 & ATACOM OK \\
0.05s  & 39.95 & 4.96$\pm$6.08 (3/5 GO) & Basin regime \\
0.10s  & 13.68 & \textbf{2.60} & DT-margin OK \\
0.20s  &  ---  & 1.33$\pm$1.50 & $h\!\gg\!h^\star$ \\
\bottomrule
\end{tabular}
\end{table}

Table~\ref{tab:dt} corroborates the M1 chord-excursion premise of the main paper's failure-mode analysis (and hence Prop.~4): continuous-time ATACOM achieves zero cost at $\Delta t \le 0.02$s but fails at $\Delta t \ge 0.05$s, the phase transition matching the M1 threshold $\Delta t\|u\| \gtrsim \sqrt{d_\text{safe}\cdot r}$. The DT-margin heuristic tracks the safe regime at $\Delta t \in \{0.10, 0.20\}$s with $h{=}3$. At $\Delta t{=}0.05$s the picture is nuanced: the horizon rule prescribes $h^\star{=}4$, but a 5-seed sweep finds $h{=}3$ at 3/5 GO (bimodal), $h{=}4$ at 0/5, and $h{=}5$ at 3/5 (one seed at 37.05). We attribute this to a \emph{BRT-shaping basin}: near the penetration horizon the lookahead reports unavoidable penetration along multiple directions every step, integrating spurious penalty into the policy gradient and trapping training.

\subsection{$\alpha$/$h$ sensitivity}
\label{sec:alpha-h-sweep}

Sweeping $\alpha$ and $h$ on Goal (5 seeds $\times$ 200K):

\begin{table}[t]
\centering
\caption{DT-margin $\alpha$/$h$ sensitivity on Goal (5 seeds $\times$ 200K).}
\label{tab:alpha-h}
\setlength{\tabcolsep}{4pt}
\begin{tabular}{lcccc}
\toprule
Config & $\alpha$ & $h$ & mean $\bar C$ & GO/5 \\
\midrule
no margin            & 0.0  & 2 & 3.37$\pm$2.81 & 4/5 \\
short horizon        & 0.05 & 1 & 1.94$\pm$1.78 & 4/5 \\
both at minimum      & 0.05 & 2 & 2.36$\pm$2.60 & 4/5 \\
above minimum        & 0.1  & 2 & \textbf{1.30$\pm$1.66} & \textbf{5/5} \\
$\alpha$ generous    & 0.3  & 1 & 1.72$\pm$1.46 & 5/5 \\
\textbf{paper choice}& 0.3  & 3 & \textbf{0.96$\pm$0.68} & \textbf{5/5} \\
over-provisioned     & 0.5  & 5 & 3.00$\pm$3.38 & 4/5 \\
\bottomrule
\end{tabular}
\end{table}

Only configurations at least one step above the minimum on \emph{both} axes ($\alpha{\ge}0.1$, $h{\ge}2$) reach 5/5 GO; setting $\alpha{=}0$ drops to 4/5 with a seed at $\bar C{=}8.46$. The over-provisioned config ($\alpha{=}0.5$, $h{=}5$) \emph{regresses} to 4/5 ($\bar C{=}9.26$ outlier) --- large-$h$ BRT shaping can integrate spurious penalty (see the $\Delta t$ sensitivity above) and destabilise training. The DT-margin recipe is thus robust only with margin on both knobs; VA-ATACOM eliminates $\alpha$ entirely, replacing it with the calibrated $a_{\max}$.

\end{document}
